\documentclass[10pt]{article}
\usepackage[a4paper,margin=2cm]{geometry}
\usepackage{booktabs}
\usepackage{array}
\usepackage{enumitem}
\setlist{nosep,leftmargin=1.2em}
\setlength{\parskip}{4pt}
\setlength{\parindent}{0pt}

\begin{document}

{\large\bfseries StepWise --- Response to Review-3 Comments}\\[2pt]
Siddhi J Salunkhe (23BDS0243), Sanvi Saket (23BDS0210)

We thank the panel for the comments. Each is answered below with what we did, what we found and
where it appears in the revised paper. Every number is regenerated by
\texttt{python -m app.ml.validate} and stored in \texttt{backend/ml\_models/reports/validation\_*.json}.
Three comments need data that only a classroom can provide. For those we report the closest
evidence available now, say plainly what is still open, and give a ready-to-run protocol
(\texttt{docs/REAL\_WORLD\_VALIDATION.md}).

\section*{Update: retraining on real handwriting and real student errors}
After the response below was drafted, both models were retrained on real data. Reports are in
\texttt{backend/ml\_models/reports/}.
\begin{itemize}
\item \textit{New real datasets.} HME100K (25,960 lines from real phone photos of students' paper
work), MathWriting 2024 (10,851 lines, many writers), MNIST train digits, and CogTutorBugs real
student errors.
\item \textit{Recognizer.} The HOG+SVM was replaced by a CNN that reads over a lattice of possible
segmentations (so touching symbols can be split). It is combined with a whole-line CTC reader
trained on 39,069 real lines and decoded with only StepWise's symbols and well-formed lines. Exact whole expressions, all test sets unseen:
\begin{center}\begin{tabular}{lrrr}\toprule
Test set & $n$ & Review-3 & Now\\\midrule
CROHME 2014/16/19 & 127 & 61.4\% & \textbf{69.3\%}\\
HME100K real phone photos & 590 & 4.6\% & \textbf{35.6\%}\\
MathWriting (unseen writers) & 34 & 41.2\% & \textbf{61.8\%}\\
Validation (1 in 7, three sources) & 697 & -- & \textbf{50.6\%}\\\bottomrule
\end{tabular}\end{center}
The target of 90\% exact expressions is not reached. On real photos most errors are touching or
crowded symbols, so confirm-before-diagnose stays. With the new reader its threshold is
$\tau=0.80$ (chosen on validation pages). Diagnosis accuracy after confirmation is 92.6\%, with
about 8 of 14 symbols checked per page. On these composed pages the new reader is exact less often
than before (9.0\% vs 16.7\%), because the line reader is not trained on pasted glyphs. A line
reading's confidence is therefore capped by the symbol CNN's, so disagreements are always
highlighted.
\item \textit{Misconception classifier (C2).} Retrained on synthetic data plus the real errors of a
30\% development split of CogTutor problem templates. On the other 70\% (1,552 real errors,
templates never seen) accuracy rises from 52.1\% to \textbf{76.6\%} (95\% CI 74.5--78.7), macro-F1
from 0.460 to \textbf{0.701}. The majority baseline is 53.2\%. All 2,120 real wrong steps are
flagged, and 99.3\% of valid steps are accepted.
\item \textit{Diverse writers (C3).} On MNIST test digits the recognizer now reads 98.8\% (high-school
students) and 99.3\% (Census employees), against 86.3\% and 85.0\% before.
\item \textit{Still open.} Real classroom study (C5) and our own phone photos (Study A). Both need
participants, and the protocols are ready in \texttt{docs/REAL\_WORLD\_VALIDATION.md}.
\end{itemize}

\section*{C1. Only 56.7\% exact expression accuracy may limit downstream diagnosis}
\textbf{Agreed, and confirmed.}
\begin{itemize}
\item \textit{Why 56.7\%.} The 127 test lines average 9.35 symbols, not about 4 as the Review-3 text
said (that figure was an error). With 92.4\% symbol accuracy, independent errors predict 54.2\%.
Lines of up to 5 symbols, the size of most student steps, are read exactly 86\% of the time. Of the
failed lines, 26 are segmentation errors (touching symbols) and 29 are substitutions.
\item \textit{Improvement.} Syntax-constrained decoding (Viterbi over the top-3 classes of each
symbol, well-formed lines only, never using arithmetic truth) raises exact lines to \textbf{61.4\%}.
\item \textit{Cost to diagnosis.} On 420 synthetic solutions written in real CROHME test-writer
handwriting, diagnosis accuracy falls from \textbf{99.3\%} on the true transcription to
\textbf{35.0\%} on the recognizer's reading, and 57\% of correct work is marked wrong.
\item \textit{Mitigation in the app.} Confirm-before-diagnose: symbols below confidence
$\tau=0.98$ (chosen on separate validation pages) are highlighted, and no diagnosis is given until the
student checks them. Accuracy returns to \textbf{93.8\%} and correct work marked wrong falls to 5\%,
at a cost of about 7 of 14 symbols checked per page (trade-off in Fig.~2). StepWise is therefore
semi-automatic until the reader improves. The study also exposed two crash bugs on misread input,
both now fixed.
\end{itemize}
\textit{Paper:} Sections IV-B, IV-C, V-B; Table VII; Fig.~2.

\section*{C2. The misconception classifier is trained on synthetic solutions; validate on real student work}
\textbf{Done with real student errors; real handwritten work still to collect.}
\begin{itemize}
\item We evaluated M2 on \textbf{CogTutorBugs} (CC BY 4.0): steps on which 95 Algebra I students
made an error caught by the Cognitive Tutor. SymPy separates 2,120 mathematically wrong steps from
3,100 valid but unhelpful ones. The data are typed, so this tests diagnosis alone.
\item \textit{Detection transfers:} 94--96\% of real wrong steps are flagged, 99.3\% of valid steps
accepted (old rule-based system: 10.8\%).
\item \textit{Naming the category does not, at first:} the Review-3 model was right for
\textbf{46.3\%} of errors on held-out problem templates, below the 53.2\% majority class (macro-F1
0.279, against 0.994 on synthetic data).
\item \textit{Why, and v2.} Error analysis on a separate 30\% of templates showed that students slip
while simplifying one side in place ($-75x=4-54 \rightarrow -75x=58$), which our generator never
produced. Nine new features and matching generators (M2 v2, still trained only on synthetic data)
give \textbf{52.1\%}, macro-F1 \textbf{0.460} on the held-out templates. In cross-validation,
adding about 160 real labelled errors lifts accuracy to \textbf{77\%}, so labelled real data is the
remaining lever.
\item \textit{Next:} Study A in \texttt{REAL\_WORLD\_VALIDATION.md}: 200 photographed pages from at
least 30 students, double-annotated (Cohen's $\kappa \ge 0.70$), evaluated by
\texttt{evaluation/real\_world\_eval.py}.
\end{itemize}
\textit{Paper:} Sections IV-E, V-E; Tables IX--X.

\section*{C3. Evaluate on actual smartphone photographs and diverse handwriting styles}
\textbf{Partly done; real phone photos are part of Study A.}
\begin{itemize}
\item \textit{Simulated capture} of the 420 pages under 12 conditions (paper and pen, notebook rules,
uneven lighting, cast shadow, tilt, blur, noise and JPEG, low resolution, three combined
severities). It exposed a design flaw: a single global threshold reads shadows and rules as ink.
A new photo-robust page binarization reduces CER on ruled paper from 0.90 to 0.16, under uneven
lighting from 0.74 to 0.10 and under a cast shadow from 0.80 to 0.12, with clean images unchanged.
Combined moderate or severe capture still defeats the reader.
\item \textit{Writer populations.} On MNIST test digits (pen on paper, never used for training) the
recognizer reads 86.3\% of digits by about 250 high-school students and 85.0\% by 48 Census
employees, against 91.0\% on CROHME digits. The main failure is a slanted `1' read as `/'.
\end{itemize}
\textit{Paper:} Sections IV-D, V-E; Table VIII.

\section*{C4. Clarify the mapping between ASSISTments skills and the eight tutoring skills}
\textbf{Done.} The new skill-mapping table gives, per StepWise skill: the ASSISTments skills, match
quality (direct: S1, S3, S4; close: S5, S6; partial: S2; approximate: S7; none: S8, which keeps
simulated parameters), data volume, the deployed BKT parameters and test AUC on those skills. We
also clarify that the AUC 0.815 benchmark uses all 110 ASSISTments skills; the mapping only
transfers BKT parameters. Two findings: several parameters sit on the EM bounds, and BKT predicts S7
barely above chance (AUC 0.54), so its S7 mastery estimate is weak.
\textit{Paper:} Section IV-F; Table XII.

\section*{C5. Validate adaptive practice through a real classroom study, not simulation alone}
\textbf{Agreed; protocol ready, trial not yet run.}
\begin{itemize}
\item \textit{Robustness now.} The deployed learner model chose items for learners from five simulators. Adaptive practice beats random with the original simulator (post-test $d=0.37$), with learners calibrated on real students ($d=0.50$) and with learners who differ in learning rate ($d=0.29$). It shrinks when learners forget ($d=0.17$) and vanishes when these are combined ($d=0.04$), because the policy stops practising skills that BKT, which has no forgetting, calls mastered. Simulation therefore cannot settle the question, as the panel noted.
\item \textit{Classroom trial (Study B):} randomized at student level, stratified by class and
pre-test, comparing adaptive with random practice from the same item bank: pre-test, four sessions
of 10 items, post-test and a two-week delayed post-test. The primary analysis is ANCOVA. Detecting
$d=0.3/0.4/0.5$ needs 132/75/49 students per arm, so a two-class first run is a pilot that estimates
the effect size and feasibility.
\end{itemize}
\textit{Paper:} Sections IV-G, V-E; Tables XIII--XIV.

\section*{Other corrections}
\begin{itemize}
\item The test-suite table now lists the real per-module counts (171 cases in 14 modules); the
Review-3 table did not add up to its total.
\item The mistake classifier is now v2 (42 features); its synthetic macro-F1 is 0.992 (was 0.994).
\item The limitations, conclusion and future work were rewritten to match the new evidence. StepWise
is presented as ready for a supervised classroom pilot, not for unsupervised use.
\end{itemize}

\end{document}
